%------------------------------------------------------------------------------%
\documentclass[a4paper,12pt,fleqn]{article}

\usepackage{calculo_varias_variaveis-1}

\ShowAnswers

%------------------------------------------------------------------------------%
\begin{document}

\makeHeader{CFVVI}{Prova 3 -- Turma 3}{20/05/2026}

\bigskip
\textbf{Atenção}: a aplicação das regras do produto ou do quociente em casos
nos quais um dos termos é constante \textbf{será considerada incorreta},
por evidenciar compreensão inadequada das regras de derivação.

% 01
%------------------------------------------------------------------------------%
\exercise{50\%}
Considere a função
\(
  f(x, y, z) = e^{xy}\ln\left(x+z^2\right) + \frac{y^2}{z} - \sin(xz)
\)
\begin{enumerate}[label=\alph*),nosep]
  \item Calcule o gradiente de $f$

  \item Calcule a derivada direcional de $f$ no ponto
  \(
    P = (1, 1, 1)
  \)
  na direção
  \(
    v = (2, -1, 2)
  \)

  \item Determine a direção em que $f$ cresce mais rapidamente no ponto $P$
\end{enumerate}

\begin{answer}
  \textbf{(a)} O gradiente de
  \(
    f(x,y,z)=e^{xy}\ln\left(x+z^2\right)+\frac{y^2}{z}-\sin(xz)
  \)
  é
  \(
    \nabla f(x, y, z)
    = \Vetortri{\D{g}{x}}{\D{g}{y}}{\D{g}{z}}
  \)
  \begin{align*}
    \D{f}{x}
    & = \D{}{x} \left( e^{xy}\ln\left(x+z^2\right) \right)
      + \D{}{x} \left( \frac{y^2}{z} \right)
      - \D{}{x} \left( \sin(xz) \right)
    \\[2mm]
    & = ye^{xy}\ln\left(x+z^2\right)
      + e^{xy}\frac{1}{x+z^2}
      - z\cos(xz)
  \end{align*}
  \begin{align*}
    \D{f}{y}
    & = \D{}{y} \left( e^{xy}\ln\left(x+z^2\right) \right)
      + \D{}{y} \left( \frac{y^2}{z} \right)
      - \D{}{y} \left( \sin(xz) \right)
    \\[2mm]
    & = xe^{xy}\ln\left(x+z^2\right)
      + \frac{2y}{z}
  \end{align*}
  \begin{align*}
    \D{f}{z}
    & = \D{}{z} \left( e^{xy}\ln\left(x+z^2\right) \right)
      + \D{}{z} \left( \frac{y^2}{z} \right)
      - \D{}{z} \left( \sin(xz) \right)
    \\[2mm]
    & = e^{xy}\frac{2z}{x+z^2}
      - \frac{y^2}{z^2}
      - x\cos(xz)
  \end{align*}
  \[
    \nabla f(x,y,z)
    = \Vetortri{
    ye^{xy}\ln(x+z^2)
    +
    \frac{e^{xy}}{x+z^2}
    -
    z\cos(xz)}{
    xe^{xy}\ln(x+z^2)
    +
    \frac{2y}{z}}{
    \frac{2ze^{xy}}{x+z^2}
    -
    \frac{y^2}{z^2}
    -
    x\cos(xz)}
  \]

  \textbf{(b)} No ponto \(P=(1, 1, 1)\),
  \[
    \nabla f(1, 1, 1)
    = \Vetortri{e\ln2+\frac{e}{2}-\cos1}{e\ln2+2}{e-1-\cos1}
  \]
  O vetor unitário na direção de
  \(
    v = (2, -1, 2)
  \)
  é
  \[
    u
    = \frac{v}{\norm{v}}
    = \frac{1}{3} \vetortri{2}{-1}{2}
    = \Vetortri{\frac{2}{3}}{-\frac{1}{3}}{\frac{2}{3}}
  \]
  A derivada direcional é
  \begin{align*}
    D_{u} f(1, 1, 1)
    & = \nabla f(1, 1, 1) \cdot u                                        \\[2mm]
    & = \Vetortri{e\ln2+\frac{e}{2}-\cos1}{e\ln2+2}{e-1-\cos1} \cdot
        \Vetortri{\frac{2}{3}}{-\frac{1}{3}}{\frac{2}{3}}                \\[2mm]
    & = \frac{2}{3} \left( e\ln 2 + \frac{e}{2} - \cos(1) \right)
      - \frac{1}{3} \left( e\ln 2 + 2 \right)
      + \frac23 \left( e - 1 - \cos(1) \right)                           \\[2mm]
    & = \frac{1}{3} \left( e\ln 2 + 3e - 4\cos(1) - 4 \right)
  \end{align*}

  \textbf{(c)} A função cresce mais rapidamente na direção do gradiente
  \[
    \nabla f(1, 1, 1)
    =\Vetortri{e\ln2+\frac{e}{2}-\cos1}{e\ln2+2}{e-1-\cos1}
  \]
\end{answer}

% 02
%------------------------------------------------------------------------------%
\exercise{50\%}
Considerando a função
\(
  g(x, y) = \sqrt{xy + 1} + \sin(x-y)
\)
\begin{enumerate}[label=\alph*),nosep]
  \item Determine a equação do plano tangente ao gráfico de $g$
        no ponto correspondente a \((0, \pi)\)
  \item Verifique se o ponto \((2, 2, -1)\) pertence ao plano
\end{enumerate}

\begin{answer}
  \textbf{(a)}

  A equação do plano tangente é

  \[
  z-z_0
  =
  g_x(x_0,y_0)(x-x_0)
  +
  g_y(x_0,y_0)(y-y_0).
  \]

  Calculamos inicialmente

  \[
  g(0,\pi)
  =
  \sqrt{1}+\sin(-\pi)
  =
  1
  \]

  Logo,

  \[
  z_0=1
  \]

  Calculamos as derivadas parciais.

  Para \(x\),

  \[
  g_x(x,y)
  =
  \frac{y}{2\sqrt{xy+1}}
  +
  \cos(x-y).
  \]

  Para \(y\),

  \[
  g_y(x,y)
  =
  \frac{x}{2\sqrt{xy+1}}
  -
  \cos(x-y).
  \]

  Avaliando em \((0,\pi)\),

  \[
  g_x(0,\pi)
  =
  \frac{\pi}{2}+\cos(-\pi)
  = \frac{\pi}{2} - 1
  \]

  \[
  g_y(0,\pi)
  =
  -\cos(-\pi)
  = 1
  \]

  \[
  z-1
  =
  \left(\frac{\pi}{2} - 1 \right)x
  + y - \pi.
  \]

  b)

  \[
    z
    =
    \left(\frac{\pi}{2} - 1 \right)2
    + y - \pi + 1
    = \pi -2 +2 -\pi + 1
    \]


\end{answer}

%------------------------------------------------------------------------------%
\end{document}
%------------------------------------------------------------------------------%

